% ===================== BASES DE DISENO =====================
\section{Bases de Diseno}
\label{sec:bases}

Las bases de dise\~{n}o del presente an\'{a}lisis corresponden a las condiciones
congeladas del proyecto P2603 SW-K60 y son comunes a las l\'{i}neas del
paquete de flexibilidad. No es admisible modificarlas sin un nuevo ciclo de
revisi\'{o}n con el cliente.

% ─────────────────────────────────────────────────────────────────────────────
\subsection{Restricciones de Operacion}
\label{ssec:restricciones}

\begin{table}[htbp]
    \centering
    \caption{Condiciones de dise\~{n}o y operaci\'{o}n de la l\'{i}nea SIM-014.}
    \label{tab:restricciones}
    \setcounter{itemcount}{0}
    \begin{tabularx}{\textwidth}{@{}NXccl@{}}
        \toprule
        \multicolumn{1}{c}{\textbf{Item}} & \textbf{Parametro} & \textbf{Valor} & \textbf{Unidad} & \textbf{Referencia} \\
        \midrule
        & Presi\'{o}n m\'{a}xima de operaci\'{o}n ($P_2$) & \num{1241} & \si{kPa}~g (180~psig) & Bases de dise\~{n}o del proyecto \\
        & Temperatura de operaci\'{o}n $T_2$ & 90 & \si{\celsius} & Bases de dise\~{n}o del proyecto \\
        & Temperatura de operaci\'{o}n alterna $T_1$ & 30 & \si{\celsius} & Modelo de an\'{a}lisis \\
        & Temperatura de instalaci\'{o}n & 25 & \si{\celsius} & Bases de dise\~{n}o del proyecto \\
        & Rango t\'{e}rmico m\'{a}ximo de evaluaci\'{o}n $\Delta T$ & 65 & \si{\celsius} & Derivado de las anteriores \\
        & C\'{o}digo de evaluaci\'{o}n & \multicolumn{2}{c}{ASME B31.3-2016} & Cliente \\
        & Criterio de aceptaci\'{o}n & \multicolumn{2}{c}{Code $\leq$ Allowable (Ratio $\leq$ 100~\%)} & ASME B31.3-2016 \\
        \bottomrule
    \end{tabularx}
\end{table}

% ─────────────────────────────────────────────────────────────────────────────
\subsection{Condiciones Ambientales del Sitio}
\label{ssec:amb}

La temperatura de instalaci\'{o}n se adopt\'{o} igual a 25~\si{\celsius}, valor
est\'{a}ndar de referencia ambiente empleado en el paquete de an\'{a}lisis; el
modelo no considera cargas de viento ni de sismo (ver limitaciones,
secci\'{o}n~\ref{sec:alcance}).

% ─────────────────────────────────────────────────────────────────────────────
\subsection{Propiedades Fisicoquimicas del Fluido de Proceso}
\label{ssec:fluido}

\begin{table}[htbp]
    \centering
    \caption{Propiedades del fluido de proceso (agua).}
    \label{tab:propfluido}
    \setcounter{itemcount}{0}
    \begin{tabularx}{\textwidth}{@{}NXccc@{}}
        \toprule
        \multicolumn{1}{c}{\textbf{Item}} & \textbf{Propiedad} & \textbf{Valor} & \textbf{Unidad} & \textbf{Referencia} \\
        \midrule
        & Fluido & \multicolumn{2}{c}{Agua de proceso} & Cliente \\
        & Gravedad espec\'{i}fica (SG) & 1.0 & --- & Bases de dise\~{n}o del proyecto \\
        & Densidad & \num{1000} & \si{kg/m^3} (0.001~\si{kg/cm^3}) & Bases de dise\~{n}o del proyecto \\
        \bottomrule
    \end{tabularx}
\end{table}

% ─────────────────────────────────────────────────────────────────────────────
\subsection{Especificacion de Tuberias y Accesorios}
\label{ssec:tuberia}

La l\'{i}nea est\'{a} construida en acero inoxidable austen\'{i}tico
ASTM A312 TP 304L seg\'{u}n ASME B36.19, Sch 10S, sin tolerancia por
corrosi\'{o}n (CA~=~0), conforme a la especificaci\'{o}n de tuber\'{i}as del
proyecto. La Tabla~\ref{tab:tuberia} resume la especificaci\'{o}n extra\'{i}da
del archivo PCF del isom\'{e}trico.

\begin{table}[htbp]
    \centering
    \caption{Especificaci\'{o}n de tuber\'{i}as y accesorios de la l\'{i}nea SIM-014 (fuente: PCF112).}
    \label{tab:tuberia}
    \setcounter{itemcount}{0}
    \begin{tabularx}{\textwidth}{@{}N>{\raggedright\arraybackslash}X>{\raggedright\arraybackslash}X@{}}
        \toprule
        \multicolumn{1}{c}{\textbf{Item}} & \textbf{Componente} & \textbf{Especificacion} \\
        \midrule
        & Tuber\'{i}a & ASTM A312 Gr TP 304L, ASME B36.19, Sch 10S, extremos PE/BW \\
        & Di\'{a}metros nominales presentes & 6 NPS; derivaciones menores de 1 NPS \\
        & Codos radio largo & ASTM A403 Gr WP 304L, ASME B16.9, Sch 10S, BW \\
        & Manguera flexible & Conexi\'{o}n de descarga en el inicio del tramo \\
        & Uniones bridadas & 150~lb, RF, ASME B16.5 \\
        & V\'{a}lvulas de compuerta de cuchilla (\textit{through conduit knife gate}) & 6 NPS, Cl.~150, tipo \textit{lug}, RF, ASME B16.10; una con actuador neum\'{a}tico \\
        & Conexiones de derivaci\'{o}n & Tee no reforzada y \textit{weldolet}, 1 NPS \\
        & Tolerancia por corrosi\'{o}n (CA) & 0 \\
        \bottomrule
    \end{tabularx}
\end{table}

% ─────────────────────────────────────────────────────────────────────────────
\subsection{Especificacion de Equipos Principales}
\label{ssec:equipos}

El equipo conectado al inicio del tramo analizado es el TMS, cuya conexi\'{o}n
de descarga, mediante manguera flexible, marca el inicio del modelo (nodo 20);
el extremo aguas abajo corresponde al nodo 170 y las derivaciones de 1~NPS
terminan en los nodos 180, 190 y 200. Las cargas que la l\'{i}nea transmite a
los soportes intermedios se reportan en la secci\'{o}n~\ref{sec:resultados}.

% ─────────────────────────────────────────────────────────────────────────────
\subsection{Normativas y Estandares Aplicables}
\label{ssec:normas}

\begin{table}[htbp]
    \centering
    \caption{Marco normativo aplicable al an\'{a}lisis de flexibilidad.}
    \label{tab:normas}
    \setcounter{itemcount}{0}
    \begin{tabularx}{\textwidth}{@{}NXXl@{}}
        \toprule
        \multicolumn{1}{c}{\textbf{Item}} & \textbf{Norma / Estandar} & \textbf{Alcance} & \textbf{Edicion} \\
        \midrule
        & ASME B31.3~\cite{asme_b313_2016} & C\'{o}digo de tuber\'{i}as de proceso: criterios de esfuerzo sostenido, expansi\'{o}n y admisibles & 2016 \\
        & ASTM A312 / ASME B36.19~\cite{astm_a312} & Material y dimensiones de tuber\'{i}a inoxidable & Vigente \\
        & ASTM A403 / ASME B16.9~\cite{asme_b169} & Accesorios forjados para soldadura a tope & Vigente \\
        & ASME B16.5~\cite{asme_b165} & Bridas y uniones bridadas & Vigente \\
        \bottomrule
    \end{tabularx}
\end{table}
